\begin{lemma}[Homogeneous quadruples]
For every infinite $X\subseteq\mathbb N$ and every two-coloring of the
four-element subsets of $X$, there are an infinite $Y\subseteq X$ and a
color $b$ such that every four-element subset of $Y$ has color $b$.
\end{lemma}
\begin{proof}
Use $F=\mathsf{CardinalityFront}(4,X)$ from
\noderef{CardinalityFront}; it is a front by
\noderef{CardinalityFrontIsFront}, and color it by
$S=\{s\in F\mid c(s)=1\}$.  Nash--Williams
\noderef{NashWilliams} supplies an infinite $Y$ and a
sub-front $F'\subseteq F$ which is either contained in $S$ or disjoint
from $S$.  Since $F'$ is nontrivial, its base is the union of its
members; because $F'\subseteq F$ and every member of $F$ is contained in
$X$, this gives $Y\subseteq X$.  The fixed cardinality argument is the same as for triples:
given any four-element $s\subseteq Y$, choose an infinite branch through
$Y$ beginning with $s$; density of $F'$ supplies a front member on that
branch, and its cardinality four forces it to be $s$.  Consequently all
four-element subsets of $Y$ receive the same color, namely $1$ in the
first case and $0$ in the second.
\end{proof}
